% Contenido del póster: Chu y Wei (2024), Journal of Financial and Quantitative Analysis. DOI 10.1017/s0022109023000698
% Elaborado con apoyo de IA a partir del resumen oficial del artículo (OpenAlex), 2026-11-12.
\def\DOI{10.1017/s0022109023000698}
\def\TituloEs{El crédito fintech y la competencia en el mercado de crédito}
\def\TituloOriginal{Fintech lending and credit market competition}
\def\Autores{Chu y Wei}
\def\Anio{2024}
\def\Revista{Journal of Financial and Quantitative Analysis}

\def\Problema{El crédito fintech crece en todo el mundo, pero no es evidente que beneficie a todos los prestatarios: sus efectos sobre el acceso, las tasas y el bienestar dependen de cómo compite con los bancos.}
\def\Objetivo{Estudiar cómo el auge del crédito fintech afecta el acceso al crédito, las tasas de interés y el bienestar social.}
\def\Metodo{Modelo de competencia crediticia con dos bancos establecidos y un prestamista fintech, que usan información y tecnologías distintas para evaluar la solvencia de los prestatarios. Los resultados se obtienen de forma analítica y numérica.}
\def\Resultados{\begin{itemize}
  \item Si la evaluación del prestamista fintech es más precisa que la de los bancos, el acceso de los prestatarios de buena calidad puede empeorar.
  \item Bajo ciertas condiciones, el crédito fintech empeora la eficiencia en la asignación del crédito y reduce el bienestar social.
  \item En la mayoría de los casos, el crédito fintech reduce las tasas de interés esperadas.
\end{itemize}}
\def\Conclusion{El crédito fintech no es beneficioso en todos los casos: aunque suele abaratar el crédito, puede perjudicar el acceso de los buenos prestatarios y la eficiencia del mercado.}
\def\Relevancia{El estudio advierte que acceso y precio pueden moverse en direcciones distintas. En el Perú, el acceso al financiamiento participativo empresarial es alto (\cifra{colocacionEmp}~\% de los proyectos publicados se financian), pero el precio no puede evaluarse porque la SMV no publica la tasa de cada operación.}
\def\PalabrasClave{crédito fintech, competencia, acceso al crédito, tasas de interés, bienestar}
% --- Ficha del estudio (herramientas/posteres/ficha_articulos.py) ---
\def\Ambito{Teórico}
\def\Enfoque{Modelo de competencia (resultados analíticos y numéricos)}
\def\FiguraPeru{2}
% --- fin de la ficha ---
